\documentclass[../../../include/open-logic-section]{subfiles}

\begin{document}

\olfileid[tr]{sth}{ord-arithmetic}{add}
\olsection{Sıra Sayısı Toplamı}

$\alpha$ ile $\beta$'yı toplamak istediğimizi varsayalım. $\alpha$'nın bir
kopyasının hemen arkasına $\beta$'nın bir {kopyasını} koyabiliriz.
(\olref[ordinals][basic]{ordinalsaresubsets} sonucundan ya
$\alpha\subseteq\beta$ ya da $\beta\subseteq\alpha$ olduğunu bildiğimiz
için 
\emph{kopyalar} almamız gerekir.) Bunun sezgisel etkisi, önce
$\alpha$ uzunluğunda bir adımlar dizisini, 
\emph{ardından} $\beta$
uzunluğunda bir adımlar dizisini katetmektir. Ortaya çıkan dizi iyi
sıralanmış olacaktır; dolayısıyla
\olref[ordinals][ordtype]{thmOrdinalRepresentation} gereği (biricik) bir
sıra sayısıyla izomorftur. Bu sıra sayısı, $\alpha$ ile $\beta$'nın
\emph{toplamı} sayılabilir.

Sıra sayısı toplamının ardındaki sezgisel fikir budur. Onu kesin biçimde
tanımlamak için kümelerin \emph{kopyalarını} alma düşüncesiyle başlarız.
Buradaki fikir, hangi nesnenin nereden geldiğini izlemek için gelişigüzel
$0$ ve $1$ etiketlerini kullanmaktır:

\begin{defn}\ollabel{defdissum}
$A$ ile $B$'nin 
\emph{ayrık toplamı}
$A\disjointsum B=(A\times\{0\})\cup(B\times\{1\})$'dir.
\end{defn}

Şimdi sıra sayısı çiftleri üzerinde bir sıralama tanımlayalım:

\begin{defn}
Her $\alpha_1,\alpha_2,\beta_1,\beta_2$ sıra sayısı için:
\begin{align*}
	\tuple{\alpha_1, \alpha_2} \rlexless \tuple{\beta_1, \beta_2}
	\text{ ancak ve ancak }& 
	\text{ya $\alpha_2 \in \beta_2$}\\
	& \text{ya da hem $\alpha_2 = \beta_2$ hem $\alpha_1 \in \beta_1$}
\end{align*}
olsun.
\end{defn}

Bu bir 
\emph{ters sözlük} sıralamasıdır; çünkü önce ikinci elemana, sonra
birinci elemana göre sıralarsınız. $\alpha\ordplus\beta$'yı,
$\alpha$'nın bir kopyasının ardından $\beta$'nın bir kopyasının geldiği
sıralamanın sıra tipi olarak tanımlamak istediğimizi hatırlayın. Bunu elde
etmek için şöyle deriz:

\begin{defn}\ollabel{defordplus}
Her $\alpha,\beta$ sıra sayısı için bunların toplamı
$\alpha\ordplus\beta=\ordtype{\alpha\disjointsum\beta,\rlexless}$'dır.
\end{defn}
\noindent
Burada gösterimi biraz kötüye kullandığımıza dikkat edin. Kesin konuşursak
``$\rlexless$'' yerine
``$\Setabs{\tuple{x,y}\in\alpha\disjointsum\beta}{x\rlexless y}$''
yazmalıydık. Yine de kısa tutmak için bundan sonra gösterimi bu biçimde
kötüye kullanmayı sürdüreceğiz.

Aşağıdaki sonuç,
\olref[ordinals][ordtype]{thmOrdinalRepresentation} ile birlikte,
tanımımızın iyi kurulmuş olduğunu doğrular:

\begin{lem}\ollabel{ordsumlessiswo}
Her $\alpha$ ve $\beta$ sıra sayısı için
$\tuple{\alpha\disjointsum\beta,\rlexless}$ bir iyi sıralamadır.
\end{lem}

\begin{proof}
$\rlexless$ bağıntısı $\alpha\disjointsum\beta$ üzerinde açıkça
bağlantılıdır. İyi temelliliği göstermek için boş olmayan bir
$X\subseteq\alpha\disjointsum\beta$ sabitleyin. $Y$, $X$'in ikinci
bileşeni olabilecek en küçük olan elemanlarının alt kümesi, yani
$Y=\Setabs{\tuple{\gamma,i}\in X}{(\forall\tuple{\delta,j}\in X)i\leq j}$
olsun. Şimdi $Y$'nin birinci bileşeni en küçük olan elemanını seçin.
\end{proof}
\noindent
Böylece sıra sayısı toplamı için güzel ve açık bir tanımımız var. Aşağıdaki
şaşırtıcı olmayan bir olgudur
(\olref[z][infinity-again]{defnomega} gereği $1=\{0\}$ olduğunu hatırlayın):
\begin{prop}
Her $\alpha$ sıra sayısı için $\alpha\ordplus 1=\ordsucc{\alpha}$.
\end{prop}

\begin{proof}
$\ordsucc{\alpha}=\alpha\cup\{\alpha\}$'dan
$\alpha\disjointsum1=(\alpha\times\{0\})\cup(\{0\}\times\{1\})$'e,
$\gamma\in\alpha$ için $f(\gamma)=\tuple{\gamma,0}$ ve
$f(\alpha)=\tuple{0,1}$ ile verilen $f$ izomorfizmasını ele alın.
\end{proof}
\noindent
Ayrıca toplamanın belirli özyineleme koşullarına uyduğunu göstermek kolaydır:

\begin{lem}\ollabel{ordadditionrecursion}
Her $\alpha,\beta$ sıra sayısı için:
\begin{align*}
	\alpha\ordplus 0 &= \alpha\\
	\alpha \ordplus (\beta\ordplus 1) &= (\alpha\ordplus\beta)\ordplus 1\\
	\alpha\ordplus\beta &= \supstrict_{\delta < \beta}(\alpha\ordplus\delta) && \text{$\beta$ bir limit sıra sayısıysa}
\end{align*}
\end{lem}

\begin{proof}
Durumları tek tek denetleyelim; önce:
\begin{align*}
	\alpha \ordplus 0
	&= \ordtype{(\alpha\times\{0\})\cup(0\times\{1\}),\rlexless}\\
	&= \ordtype{\alpha\times\{0\},\rlexless}\\
	&= \alpha\\
	\alpha \ordplus (\beta\ordplus 1)
	%&= \ordtype{(\alpha\times \{0\}) \cup (\ordtype{\beta \ordplus  1}\times \{1\}), \rlexless} \\
	&= \ordtype{(\alpha\times\{0\})\cup(\ordsucc{\beta}\times\{1\}),\rlexless}\\
	&= \ordtype{(\alpha\times\{0\})\cup(\beta\times\{1\}),\rlexless}\ordplus 1\\
	&= (\alpha\ordplus\beta)\ordplus 1
\end{align*}
Şimdi $\beta\neq\emptyset$ bir limit sıra sayısı olsun. $\delta<\beta$ ise
$\delta\ordplus 1<\beta$ da olur; dolayısıyla
$\alpha\ordplus\delta$, $\alpha\ordplus\beta$'nın öz başlangıç kesitidir.
Bu nedenle $\alpha\ordplus\beta$,
$X=\Setabs{\alpha\ordplus\delta}{\delta<\beta}$ için sıkı bir üst
sınırdır. Ayrıca $\alpha\leq\gamma<\alpha\ordplus\beta$ ise açıkça bazı
$\delta<\beta$ için $\gamma=\alpha\ordplus\delta$ olur. Dolayısıyla
$\alpha\ordplus\beta=\supstrict_{\delta<\beta}(\alpha\ordplus\delta)$.
\end{proof}

Fakat çarpıcı bir olgu şudur: Sıra sayısı toplamını tanımlamak için bunun
\emph{yerine} yalnızca Sonlu Ötesi Özyineleme Teoremi'ni kullanabilir ve
özyineleme denklemlerini \olref{ordadditionrecursion} içinde verildiği gibi
(yalnız ``$\beta\ordplus 1$'' yerine ``$\ordsucc{\beta}$'' kullanarak)
koyabilirdik.

Öyleyse sıra sayıları üzerindeki işlemleri tanımlamanın iki farklı yolu
vardır. Açıkça iyi sıralanmış bir küme kurup onun sıra tipini ele alarak
\emph{sentetik} bir tanım verebiliriz. Ya da yalnızca özyineleme denklemlerini
koyarak 
\emph{özyinelemeli} bir tanım verebiliriz. Bunlar doğru yapıldığında
sonuç aynıdır. Nitekim
\olref[ordinals][ordtype]{thmOrdinalRepresentation}, bu özyineleme
denklemlerinin 
\emph{biricik} sıra sayılarını belirlemesini güvenceye alır.

Sıra sayısı aritmetiği birçok bakımdan doğal sayıların toplamı gibi davranır.
Örneğin aşağıdakini ispatlayabiliriz:

\begin{lem}\ollabel{ordinaladditionisnice}
$\alpha,\beta,\gamma$ sıra sayılarıysa:
\begin{enumerate}
	\item\ollabel{ordaddition1} $\beta<\gamma$ ise
	$\alpha\ordplus\beta<\alpha\ordplus\gamma$
	\item\ollabel{ordaddition2} $\alpha\ordplus\beta=
	\alpha\ordplus\gamma$ ise $\beta=\gamma$
	\item\ollabel{ordaddition3} $\alpha\ordplus(\beta\ordplus\gamma)
	=(\alpha\ordplus\beta)\ordplus\gamma$, yani toplama birleşmelidir
	\item\ollabel{ordaddition4} $\alpha\leq\beta$ ise
	$\alpha\ordplus\gamma\leq\beta\ordplus\gamma$
\end{enumerate}
\end{lem}

\begin{proof}
Geri kalanını alıştırma olarak bırakıp \olref{ordaddition3}'ü ispatlayalım.
İspat, \olref{ordadditionrecursion} kullanılarak $\gamma$ üzerinde Basit
Sonlu Ötesi Tümevarım'ladır. $\gamma=0$ olduğunda:
\[
(\alpha\ordplus\beta)\ordplus0=\alpha\ordplus\beta
=\alpha\ordplus(\beta\ordplus0)
\]
$\gamma=\delta\ordplus1$ olduğunda tümevarım varsayımı olarak
$(\alpha\ordplus\beta)\ordplus\delta
=\alpha\ordplus(\beta\ordplus\delta)$ alalım. Şimdi
\olref{ordadditionrecursion} sonucunu üç kez kullanarak:
\begin{align*}
	(\alpha \ordplus  \beta) \ordplus  (\delta \ordplus  1) & = ((\alpha \ordplus  \beta) \ordplus  \delta)\ordplus 1\\
	& = (\alpha \ordplus  (\beta \ordplus  \delta)) \ordplus  1\\
	& = \alpha \ordplus  ((\beta \ordplus  \delta)\ordplus 1)\\
	& = \alpha \ordplus  (\beta \ordplus  (\delta\ordplus 1))
\end{align*}
$\gamma$ bir limit sıra sayısı olduğunda, $\delta\in\gamma$ ise
$(\alpha\ordplus\beta)\ordplus\delta
=\alpha\ordplus(\beta\ordplus\delta)$ olduğunu tümevarım varsayımı olarak
alalım. Şimdi:
\begin{align*}
	(\alpha \ordplus  \beta) \ordplus  \gamma & = \supstrict_{\delta < \gamma}((\alpha \ordplus  \beta) \ordplus  \delta) \\
	&= \supstrict_{\delta < \gamma}(\alpha \ordplus  (\beta \ordplus  \delta))\\
	&= \alpha \ordplus  \supstrict_{\delta < \gamma}(\beta \ordplus  \delta)\\
	& = \alpha \ordplus  (\beta \ordplus  \gamma)
\end{align*}
\end{proof}

\begin{prob}
\olref[sth][ord-arithmetic][add]{ordinaladditionisnice} sonucunun geri
kalanını ispatlayın.
\end{prob}

Bu bakımlardan sıra sayısı toplamı son derece tanıdık gelmelidir. Fakat sıra
sayısı toplamının doğal sayılar üzerindeki toplamaya 
\emph{benzemediği} çok
önemli bir yön vardır.

\begin{prop}\ollabel{ordsumnotcommute}
Sıra sayısı toplamı {değişmeli} değildir;
$1\ordplus\omega=\omega<\omega\ordplus1$.
\end{prop}

\begin{proof}
$1\ordplus\omega=\supstrict_{n<\omega}(1\ordplus n)
=\omega\in\omega\cup\{\omega\}=\ordsucc{\omega}=\omega\ordplus1$
olduğuna dikkat edin.
\end{proof}
\noindent
Bu ilk bakışta şaşırtıcı gelebilir, fakat 
\emph{gelmemelidir}. Bir yandan
$1\ordplus\omega$'yı ele aldığınızda bütün doğal sayıların
\emph{önüne} fazladan bir eleman koyarak elde edilen sıra tipini
düşünürsünüz. \olref[sfr][infinite][hilbert]{sec} içindeki Hilbert Oteli
örneğinde yürüttüğümüz akıl yürütmeye göre bu fazladan ilk eleman, sezgisel
olarak toplam sıra tipini değiştirmemelidir. Öte yandan
$\omega\ordplus1$'i ele aldığınızda bütün doğal sayıların 
\emph{arkasına}
fazladan bir eleman koyarak elde edilen sıra tipini düşünürsünüz. Bu ise
bütünüyle farklı bir varlıktır!

\end{document}
